% =============================================
% CHƯƠNG 1 — GIỚI THIỆU ĐỀ TÀI (6 trang)
% Status: OUTLINE — chờ draft từ agent
% =============================================

\chapter{Giới thiệu đề tài}
\label{ch:gioi-thieu}

% [INTRO CHƯƠNG — 1 đoạn ~5 dòng]
% Giới thiệu chương sẽ nói về: bối cảnh điểm danh truyền thống, vấn đề thực tiễn,
% mục tiêu của đồ án, phạm vi giải quyết, đóng góp chính, và cấu trúc các chương tiếp theo.

% ─────────────────────────────────────────────
\section{Bối cảnh}
\label{sec:boi-canh}
% Ước lượng: 1 trang
% Nguồn tài sản:
%   - báo cáo tuần 16-21/03 mục "Thông tin chung"
%   - CLAUDE.md mục Project Overview
%   - Notion page 1
% Nội dung cần có:
%   - Thực trạng điểm danh tại trường ĐH (gọi tên, ký giấy, BLE Beacon, RFID, QR tĩnh...)
%   - Hiện trạng tại ĐHBKHN cụ thể (nếu có data)
%   - Bùng nổ điện thoại thông minh + Zalo (~75M user VN) tạo cơ hội
%   - Tham chiếu IEEE 1-2 paper về smart attendance systems

% ─────────────────────────────────────────────
\section{Vấn đề thực tiễn và động lực}
\label{sec:van-de}
% Ước lượng: 1 trang
% Nguồn:
%   - báo cáo 11-16/05 mục "Bối cảnh & vấn đề" (bottleneck QR nhỏ trên điện thoại)
%   - báo cáo 04-09/05 (8 UX issues khi pilot)
% Nội dung:
%   - Vấn đề 1: gian lận điểm danh hộ (proxy attendance)
%   - Vấn đề 2: tốc độ chậm khi lớp đông >50 SV
%   - Vấn đề 3: GV khó giám sát realtime, không có báo cáo dài hạn
%   - Vấn đề 4: SV không có cách xin nghỉ chính thức trong hệ thống
%   - Đặt câu hỏi nghiên cứu: làm thế nào kết hợp QR động + sinh trắc + xác minh ngang hàng
%     để vừa nhanh vừa chống gian lận?

% ─────────────────────────────────────────────
\section{Mục tiêu của đồ án}
\label{sec:muc-tieu}
% Ước lượng: 1 trang
% VIẾT MỚI hoàn toàn
% Nội dung:
%   - Mục tiêu tổng quát (1 câu)
%   - 3-5 mục tiêu cụ thể, đo lường được:
%     1. Xây dựng Mini App điểm danh dùng QR động 30s + HMAC-SHA256
%     2. Tích hợp xác minh khuôn mặt (face verification) với ngưỡng ≥70%
%     3. Cơ chế xác minh ngang hàng (peer-to-peer) tối thiểu 3 sinh viên
%     4. Dashboard quản trị web cho phòng đào tạo + giảng viên
%     5. Hệ thống chống gian lận 5 lớp với báo cáo tự động
%   - Bảng tiêu chí thành công (SMART)

% ─────────────────────────────────────────────
\section{Phạm vi đề tài}
\label{sec:pham-vi}
% Ước lượng: 1 trang
% Nguồn: memory project-constraints.md
% Nội dung:
%   - Phạm vi nghiệp vụ: SV - GV - Admin (3 vai trò)
%   - Phạm vi kỹ thuật:
%     * Frontend: Zalo Mini App (React 18 + TypeScript + Jotai)
%     * Web Admin: Vite + Ant Design
%     * Backend: Firebase Firestore + Cloud Functions (Spark plan)
%   - Phạm vi triển khai: pilot trong khuôn khổ đồ án, KHÔNG triển khai production toàn trường
%   - Phạm vi loại trừ: không tích hợp LMS, không tự động sinh báo cáo PDF cho phòng đào tạo,
%     không tích hợp thẻ sinh viên RFID

% ─────────────────────────────────────────────
\section{Đóng góp chính}
\label{sec:dong-gop}
% Ước lượng: 1 trang
% VIẾT MỚI hoàn toàn
% Nội dung:
%   - Đóng góp 1: Thiết kế cơ chế QR HMAC-SHA256 xoay 30s với phối hợp wall-clock đồng bộ
%     giữa nhiều thiết bị (chi tiết Ch.4)
%   - Đóng góp 2: Cơ chế "Cổng máy chiếu" (Teacher Display Pairing) qua QR token 128-bit,
%     one-shot transition (chi tiết Ch.4)
%   - Đóng góp 3: Thuật toán phát hiện gian lận 5 lớp (multi-tier anti-fraud)
%   - Đóng góp 4: Hệ thống 3 vai trò với phân quyền Firestore rules + Custom Token
%   - Đóng góp 5: Tích hợp tuân thủ chính sách Zalo (privacy, AI proxy, consent)

% ─────────────────────────────────────────────
\section{Bố cục quyển báo cáo}
\label{sec:bo-cuc}
% Ước lượng: 1 trang
% VIẾT MỚI
% Nội dung: liệt kê 5 chương + Mở đầu + Kết luận + Phụ lục, mỗi chương 2-3 dòng tóm tắt

% ─────────────────────────────────────────────
% \section*{Kết chương}
% 1 đoạn 3-4 dòng tóm tắt chương 1, dẫn dắt sang chương 2 (khảo sát).
